% =====================================================================
% Supplementary Document for the Deep Double Descent Poster
% COMP34312 Coursework, Group: Ayush Kumar, Jacob Yip, Harshita Dassani
%
% Purpose: backs up every claim on the poster.
%   1. Every equation derived or sourced.
%   2. Every reference annotated with what we use from it.
%   3. Critical-thinking reasoning chains for likely Q&A.
%
% NOT for submission. For internal group review and Q&A rehearsal.
% Compile: pdflatex supplementary.tex
% =====================================================================
\documentclass[11pt,a4paper]{article}

\usepackage[margin=2.2cm]{geometry}
\usepackage{amsmath,amssymb,amsthm,bm,mathtools}
\usepackage{xcolor}
\usepackage[hidelinks]{hyperref}
\usepackage{enumitem}
\usepackage{microtype}
\usepackage{booktabs}
\usepackage{titlesec}
\usepackage{lmodern}

\titleformat{\section}{\large\bfseries\color{black!85}}{\thesection.}{0.5em}{}
\titleformat{\subsection}{\normalsize\bfseries\color{black!75}}{\thesubsection}{0.5em}{}

\definecolor{uompurple}{HTML}{660099}
\newcommand{\kw}[1]{\textcolor{uompurple}{\textbf{#1}}}

\newcommand{\E}{\mathbb{E}}
\newcommand{\T}{\mathcal{T}}
\newcommand{\D}{\mathcal{D}}
\newcommand{\F}{\mathcal{F}}
\newcommand{\Sset}{\mathcal{S}}
\newcommand{\X}{\mathcal{X}}
\newcommand{\Y}{\mathcal{Y}}
\newcommand{\Hyp}{\mathcal{H}}

\theoremstyle{definition}
\newtheorem{defn}{Definition}
\newtheorem{claim}{Claim}

\title{Supplementary Backup for the\\``Deep Double Descent'' Poster}
\author{Ayush Kumar \quad Jacob Yip \quad Harshita Dassani \\[2pt]
\small COMP34312 Mathematical Topics in Machine Learning, 2025--26}
\date{}

\begin{document}
\maketitle
\vspace{-2.5em}

\begin{center}
\textit{Internal study aid. Backs up every equation, reference, and claim on the poster, plus implication arguments for Q\&A.}
\end{center}

\hrule
\vspace{0.6em}

\tableofcontents
\vspace{0.4em}
\hrule

% =====================================================================
\section{Coursework proof requirement}
% =====================================================================

The coursework specification says that when a theorem or lemma number is selected as a suggested topic, its proof is expected to be part of the presentation. It also says that if a proof is presented, proofs of invoked results should also be presentable.

Our poster does \emph{not} present a new theorem proof. It uses definitions, an algebraic identity from the notes, empirical findings from Nakkiran et al., and one informal hypothesis. Therefore the working standard for the group is:
\begin{itemize}[leftmargin=1.2em]
\item Every displayed equation must be explainable without notes.
\item Each mathematical statement must be classified as a definition, identity, empirical observation, cited theorem, or hypothesis.
\item We should not present Nakkiran et al.'s deep-net claim as a theorem. It is an empirically supported hypothesis.
\item If we choose to mention a rigorous theorem, such as Mei \& Montanari's random-features result, we should state only the high-level content unless we are ready to discuss its assumptions.
\end{itemize}

% =====================================================================
\section{Equation index}
% =====================================================================

Every equation on the poster, with its source and (where applicable) a one-step derivation. Notation matches the COMP34312 lecture notes.

\subsection{Population risk \texorpdfstring{$R(f)$}{R(f)}}

\noindent\textbf{Source.} Notes \S1.3, p.~11.

\begin{defn}[Population risk]
For data distribution $\D$ on $\X\times\Y$, a model $f:\X\to\Y$, and a non-negative loss $\ell:\Y\times\Y\to\mathbb{R}_{\ge 0}$,
\[
R(f) := \E_{(\bm{x},\bm{y})\sim\D}\bigl[\ell(\bm{y},f(\bm{x}))\bigr].
\]
\end{defn}

\noindent\textbf{Why we use it.} It is the quantity any honest model wants to minimise but cannot directly compute (we don't know $\D$). Test error is a Monte-Carlo estimate of $R(f)$.

\subsection{Empirical risk \texorpdfstring{$\widehat R(f,\Sset)$}{R-hat(f, S)}}

\noindent\textbf{Source.} Notes \S1.3, p.~11.

\begin{defn}[Empirical risk]
For a sample $\Sset = \{(\bm{x}_i,\bm{y}_i)\}_{i=1}^{n}\sim\D^n$,
\[
\widehat R(f,\Sset) := \frac{1}{n}\sum_{i=1}^{n}\ell(\bm{y}_i,f(\bm{x}_i)).
\]
\end{defn}

\noindent\textbf{Relation to $R(f)$.} For a \emph{fixed} $f$ independent of $\Sset$, $\E_{\Sset\sim\D^n}[\widehat R(f,\Sset)] = R(f)$ by linearity of expectation. The danger is that $f = \hat f$ depends on $\Sset$, so this unbiasedness fails.

\subsection{Trained model \texorpdfstring{$\hat f$}{f-hat} and the four-model hierarchy}

\noindent\textbf{Source.} Notes Defs 11--14, pp.~47.

\begin{align*}
\hat f &\leftarrow \mathcal{A}(\F;\Sset) && \text{(trained model: algorithm output on a finite sample)}\\
f_{\mathrm{erm}} &:= \arg\min_{f\in\F}\widehat R(f,\Sset) && \text{(empirical risk minimiser in $\F$)}\\
f^{*} &:= \arg\min_{f\in\F}R(f) && \text{(best-in-class with infinite data)}\\
y^{*} &:= \arg\min_{f\in\Omega}R(f) && \text{(Bayes-optimal predictor in $\Omega\supseteq\F$)}
\end{align*}

\noindent\textbf{Increasing assumptions.} $\hat f \to f_{\mathrm{erm}}$ assumes a \emph{perfect optimiser}; $f_{\mathrm{erm}}\to f^{*}$ assumes \emph{infinite data}; $f^{*}\to y^{*}$ assumes the function class is unrestricted.

\subsection{Excess-risk decomposition}

\noindent\textbf{Source.} Notes \S5.1, Eq.~(73), p.~48.

\begin{align}
R(\hat f) - R(y^{*}) &= \bigl[R(\hat f) - R(f_{\mathrm{erm}})\bigr] + \bigl[R(f_{\mathrm{erm}}) - R(f^{*})\bigr] + \bigl[R(f^{*}) - R(y^{*})\bigr]\nonumber\\
&= \underbrace{\text{optimisation}}_{\ge 0} + \underbrace{\text{estimation}}_{\ge 0} + \underbrace{\text{approximation}}_{\ge 0}. \label{eq:excess}
\end{align}

\noindent\textbf{Derivation.} Pure algebra: add and subtract $R(f_{\mathrm{erm}})$ and $R(f^{*})$ inside $R(\hat f) - R(y^{*})$. Each bracket is $\ge 0$ because each model in the chain $\hat f \to f_{\mathrm{erm}} \to f^{*} \to y^{*}$ is a worse minimiser than the next.

\subsection{Bias--variance to approximation/estimation/optimisation bridge}

\noindent\textbf{Source.} Notes \S5.1, p.~50, plus \S3.3 for the bias--variance identity.

For squared loss with a fixed target, the bias--variance identity over random samples $\Sset$ is
\[
\E_{\Sset}\bigl[(y - \hat f(\bm{x}))^{2}\bigr] = \underbrace{\bigl(\E_{\Sset}[\hat f(\bm{x})] - y\bigr)^{2}}_{\text{bias}^2} + \underbrace{\E_{\Sset}\bigl[(\hat f(\bm{x}) - \E_{\Sset}[\hat f(\bm{x})])^{2}\bigr]}_{\text{variance}} + \sigma^{2}_{\text{noise}}.
\]

\noindent\textbf{Informal mapping to (\ref{eq:excess}).}
\begin{align*}
\text{bias}^{2}\;\;&\approx\;\;\text{approximation} + (\text{part of}) \text{ estimation},\\
\text{variance}\;\;&\approx\;\;(\text{rest of}) \text{ estimation} + \text{ optimisation}.
\end{align*}

\noindent\textbf{Key point.} The bias--variance identity is a \emph{decomposition of squared error}, not a refutation of double descent. What changes in the modern regime is that ``variance'' (informally, sensitivity of $\hat f$ to $\Sset$) is no longer monotone in raw model size, because $\T$ controls which interpolant the optimiser picks.

\subsection{Effective Model Complexity (EMC), Def.~1 of Nakkiran et al.~(2019)}

\noindent\textbf{Source.} Nakkiran et al.\ (2019) \S2, Def.~1.

\begin{defn}[Effective Model Complexity]
A \kw{training procedure} $\T$ takes a labelled sample $\Sset = \{(\bm{x}_i,\bm{y}_i)\}_{i=1}^n$ and returns a classifier $\T(\Sset)$. For an error tolerance $\varepsilon>0$, the EMC of $\T$ on $\D$ is
\[
\mathrm{EMC}_{\D,\varepsilon}(\T) \;:=\; \max\bigl\{\,n\;\bigm|\;\E_{\Sset\sim\D^{n}}\bigl[\mathrm{Error}_{\Sset}(\T(\Sset))\bigr]\le\varepsilon\,\bigr\},
\]
where $\mathrm{Error}_{\Sset}(M) = \widehat R(M, \Sset)$ is the empirical error of model $M$ on $\Sset$.
\end{defn}

\noindent\textbf{Three things to internalise.}
\begin{itemize}
\item $\T$ is the \emph{whole training procedure}: architecture, optimiser, learning rate, epochs, augmentation, regularisation, label policy. Two procedures with the same architecture can have different EMCs.
\item The expectation is over training samples, not test samples. EMC asks: \emph{how much data can the procedure absorb?}
\item $\varepsilon$ is heuristic. The paper uses $\varepsilon = 0.1$. The qualitative phenomenon is robust to small perturbations of $\varepsilon$ but the critical-interval \emph{width} depends on it.
\end{itemize}

\subsection{Generalised Double Descent Hypothesis (Hyp.~1 of Nakkiran et al.)}

\noindent\textbf{Source.} Nakkiran et al.\ (2019) \S2, Hypothesis 1 (informal).

For ``natural'' $\D$, neural-network procedures $\T$, and small $\varepsilon$:
\begin{itemize}
\item If $\mathrm{EMC}_{\D,\varepsilon}(\T) \ll n$ (under-parameterised): any perturbation of $\T$ that increases EMC \emph{decreases} test error.
\item If $\mathrm{EMC}_{\D,\varepsilon}(\T) \approx n$ (critically parameterised): the same perturbation \emph{may decrease or increase} test error.
\item If $\mathrm{EMC}_{\D,\varepsilon}(\T) \gg n$ (over-parameterised): any perturbation that increases EMC \emph{decreases} test error.
\end{itemize}

\noindent\textbf{What ``hypothesis'' means here.} It is an empirical pattern that the paper validates across architectures and tasks, supplemented by linear-model analogues. It is \emph{not} a theorem about deep nets.

\subsection{Linear-model rigorous result we cite}

\noindent\textbf{Source.} Mei \& Montanari, \emph{Comm.\ Pure Appl.\ Math.} (2022).

For random-features regression on $\mathbb{R}^d$ with $N$ features, $n$ samples, ridge $\lambda\to 0^+$, and an isotropic Gaussian model, the test risk has a peak at $N = n$. Formally, in the proportional asymptotic regime $n,d,N \to\infty$ with $N/n\to\psi_1$ and $d/n\to\psi_2$, the test risk admits a closed form (their Thm.~2) that diverges as $\psi_1\to 1$ when $\lambda\to 0$.

This is the \emph{rigorous} backbone of the deep-net intuition: ``one interpolant is fragile; many interpolants give the optimiser room.''

% =====================================================================
\section{Reference index}
% =====================================================================

Every reference on the poster, with what we actually take from it.

\renewcommand{\arraystretch}{1.18}
\noindent\begin{tabular}{@{}p{4.4cm}p{4cm}p{6.5cm}@{}}
\toprule
Reference & Type & What we use from it \\
\midrule
Nakkiran et al.\ (2019), arXiv:1912.02292 & primary paper & Def.~1 (EMC), Hyp.~1, Figs.~1--3, mechanism intuition. \\
Belkin, Hsu, Ma, Mandal (2019), \emph{PNAS} & precursor & First framing of model-wise double descent for general predictors; cited as the conceptual ancestor of Nakkiran. \\
Zhang, Bengio, Hardt, Recht, Vinyals (2017), ICLR & precursor & Random-label experiments showing deep nets can memorise; motivates ``capacity $\ne$ parameter count.''  \\
Mei \& Montanari (2022), \emph{CPAM} & rigorous backing & Closed-form test risk for random-features regression with peak at $N=n$. \\
Bartlett, Long, Lugosi, Tsigler (2020), \emph{PNAS} & rigorous backing & Benign overfitting under min-norm interpolation in linear regression. \\
Hastie, Tibshirani, Friedman (2009) & textbook & Classical bias--variance trade-off, U-curve. \\
Krizhevsky et al.\ (2012); Huang et al.\ (2018) & motivation & Empirical ``larger models are better'' folklore that the paper interrogates. \\
Brown \& Mukherjee, COMP34312 Notes 2026 & course & \S5.1 excess-risk decomposition (Eq.~73); \S5.3 introduction to double descent; notation page. \\
\bottomrule
\end{tabular}

\medskip

\noindent\textbf{Notes on how we cite.}
\begin{itemize}[leftmargin=1.2em]
\item We never cite a theorem we cannot state in one sentence under questioning.
\item We never cite a paper as ``proving'' something it only conjectures (this is the trap with neural network papers in this area).
\item Belkin and Mei \& Montanari are cited together because Belkin frames the phenomenon and Mei \& Montanari prove a tractable instance.
\end{itemize}

% =====================================================================
\section{Critical-thinking reasoning chains}
% =====================================================================

Each chain takes a hypothetical perturbation and walks through its consequences using EMC and the hypothesis. These are the kinds of questions the marker is likely to ask, and answering them with a reasoned argument rather than memorised bullets is the rubric line 1b move.

\subsection{Chain 1: ``What if we had $n/2$ training samples?''}

\textbf{Setup.} Hold $\T$ fixed. Halve the sample count from $n$ to $n/2$.

\textbf{Reasoning chain.}
\begin{enumerate}[leftmargin=1.4em]
\item EMC$_{\D,\varepsilon}(\T)$ is a property of $\T$ and $\D$, \emph{independent of the actual sample count we use}, so EMC is unchanged.
\item What changes is the ratio $\mathrm{EMC}/n$, which doubles.
\item If we were originally in the over-parameterised regime ($\mathrm{EMC} \gg n$), we are now \emph{further} into it. The sample-wise double-descent argument predicts a possible local improvement, but this competes with the ordinary cost of using fewer examples.
\item If we were originally in the critical regime ($\mathrm{EMC}\approx n$), we may now move into the over-parameterised side. Test error could improve.
\item If we were already in the under-parameterised regime ($\mathrm{EMC}\ll n$), halving $n$ probably makes things worse because variance grows like $1/n$.
\end{enumerate}

\textbf{Headline.} ``More data hurts'' is local; ``less data hurts'' is asymptotic. The two facts coexist because they refer to different regimes.

\subsection{Chain 2: ``What if we trained for half as many epochs?''}

\textbf{Setup.} Hold architecture, $\F$, and $n$ fixed. Define $\T_t$ = $\T$ with $t$ epochs. Move from $t$ to $t/2$.

\textbf{Reasoning chain.}
\begin{enumerate}[leftmargin=1.4em]
\item In the paper's framing and experiments, EMC generally increases with training time because longer training gives the optimiser more opportunity to fit the sample. Strict monotonicity is a heuristic, not a theorem.
\item Halving $t$ therefore usually \emph{decreases} EMC.
\item Whether this helps or hurts depends on which side of $n$ we sit on.
\item If we were over-parameterised, decreasing EMC moves us toward the critical regime, which is bad.
\item If we were critically parameterised, decreasing EMC moves us into the under-parameterised regime, so the effect is ambiguous and test error could go either way.
\item If we were under-parameterised, decreasing EMC further is just under-fitting, which is bad.
\end{enumerate}

\textbf{Headline.} Early stopping helps narrowly, in the band of ``medium-sized'' models that would otherwise pass through the epoch-wise peak.

\subsection{Chain 3: ``What if we increased the label noise from 15\% to 30\%?''}

\textbf{Reasoning chain.}
\begin{enumerate}[leftmargin=1.4em]
\item Label noise raises the entropy of the conditional $\D(\bm{y}\mid\bm{x})$. Achieving train error below $\varepsilon$ requires the model to memorise more idiosyncratic patterns.
\item Effectively, the threshold ``how much data can $\T$ fit to within $\varepsilon$?'' falls. EMC drops.
\item For the same $n$, falling EMC moves the system from over-parameterised toward critical.
\item Additionally, near the critical regime, fitting more noise destroys global structure more aggressively. The peak height grows even if its location is the same.
\end{enumerate}

\textbf{Headline.} Label noise both \emph{shifts} the peak and \emph{amplifies} it, but does not \emph{cause} it (Fig.\ 4a of the paper shows the peak at 0\% noise too).

\subsection{Chain 4: ``What if we added strong weight decay?''}

\textbf{Reasoning chain.}
\begin{enumerate}[leftmargin=1.4em]
\item Weight decay shrinks the set of solutions $\T$ explores, so only low-norm interpolants are reachable.
\item This \emph{reduces} EMC (some sample sizes that were fittable without decay are no longer fittable with it).
\item In the over-parameterised regime, lower EMC at fixed $n$ moves us toward the critical regime, which is bad on its own.
\item But weight decay also \emph{biases the optimiser} toward the same low-norm interpolant the over-parameterised regime relies on. So it can suppress the peak entirely.
\end{enumerate}

\textbf{Headline.} Weight decay does two things at once. The implicit-bias effect dominates empirically, which is why regularisation flattens double descent in practice (paper \S5).

\subsection{Chain 5: ``What if we ensembled $K$ independent models?''}

\textbf{Reasoning chain.}
\begin{enumerate}[leftmargin=1.4em]
\item Each individual model $\hat f_k$ is a noisy interpolant near the critical regime.
\item Their average $\bar f = (1/K)\sum_k \hat f_k$ keeps the on-distribution prediction (where individual models agree) and cancels the high-variance off-distribution noise (where they disagree because they fit noise differently).
\item This is exactly the bias--variance identity: variance of the average drops as $1/K$ if the models are uncorrelated.
\item In double-descent terms: ensembling acts like an extra capacity boost that selects a better interpolant. Paper Fig.\ 28 shows ensembling flattens the peak.
\end{enumerate}

\textbf{Headline.} Ensembling is a clean experimental knockout argument that the peak is about \emph{which} interpolant gets picked, not about interpolation per se.

\subsection{Chain 6: ``What if SGD is replaced by full-batch GD?''}

\textbf{Reasoning chain.}
\begin{enumerate}[leftmargin=1.4em]
\item Full-batch GD has different implicit bias: it minimises training loss exactly along the gradient flow, with no stochastic averaging.
\item In over-parameterised settings, full-batch GD on least-squares converges to the \emph{minimum-norm} solution among interpolants. This is the same as SGD asymptotically for linear models, but with less stochastic regularisation along the way.
\item For deep nets, the implicit bias is poorly understood for both SGD and full-batch GD. Empirically, both show double descent (paper Fig.\ 6 compares).
\item So Hypothesis 1 is robust to the SGD/GD distinction in practice, but the \emph{theoretical} mechanism (which interpolant gets picked) is more transparent for one than the other.
\end{enumerate}

\textbf{Headline.} The critical regime is defined by EMC vs $n$, not by the optimiser. The optimiser changes the height/width of the peak but not its existence.

\subsection{Chain 7: ``What if the data distribution itself were noiseless and well-specified?''}

\textbf{Reasoning chain.}
\begin{enumerate}[leftmargin=1.4em]
\item In the noiseless, perfectly specified case, there exists $f^{*}\in\F$ with $R(f^{*})=0$, so the approximation gap vanishes.
\item The estimation gap is also driven down: every interpolant agrees with the true labels, so the choice among interpolants matters less.
\item The optimisation gap can still produce some sensitivity, but it shrinks too.
\item Prediction: the peak should \emph{shrink} but not necessarily disappear. Some peak survives if $\F$ is mis-specified relative to $\D$ even without label noise (paper Fig.\ 7 confirms this).
\end{enumerate}

\textbf{Headline.} Double descent is fundamentally about \emph{model mis-specification}, not just noisy labels. Eliminating noise narrows the peak; eliminating mis-specification too should remove it almost entirely.

\subsection{Chain 8: ``What if we used a much larger architecture (width $> 100$)?''}

\textbf{Reasoning chain.}
\begin{enumerate}[leftmargin=1.4em]
\item Increasing width generally raises EMC because the model class and the set of reachable fits become larger. For trained deep nets this is an empirical tendency, not a formal monotonicity theorem.
\item For the heatmap experiment ($n = 50000$ on CIFAR-10 with 15\% label noise), EMC at width 64 is already $\gg n$ for sufficient training time.
\item Going further past 64 just deepens the over-parameterised regime: test error continues to improve, but the marginal returns diminish (paper Fig.\ 4: test error plateaus past width 30 in the noiseless case).
\item No second peak is expected because there is no second critical regime to cross.
\end{enumerate}

\textbf{Headline.} Once you are clearly beyond the threshold, scale gives you returns but not surprises. The drama is local to $\mathrm{EMC}\approx n$.

% =====================================================================
\section{Anti-claims with full justification}
% =====================================================================

The poster has a ``What we do not claim'' block. Each anti-claim is one bullet on the poster; here is the longer-form justification you would give if pressed.

\subsection{``Bias--variance is wrong.'' \textbf{No.}}

The bias--variance identity (\S1.5) is an algebraic identity for squared loss. It is a \emph{decomposition}, not a predictive theory. What is incomplete is the \emph{intuition} that came with it: that variance is monotone in raw model size and overwhelms test error past a threshold. That intuition assumes the empirical risk minimiser is unique; in over-parameterised settings, it is not, and the optimiser's choice among interpolants matters.

\subsection{``Bigger / longer / more data is always better.'' \textbf{No.}}

Each of these can be locally worse inside the critical interval. Hypothesis 1 says: the sign of the change depends on which regime you sit in. The asymptotic ``more is better'' statements remain true; the \emph{local non-monotonicity} is the new content.

\subsection{``Label noise causes the peak.'' \textbf{Partially.}}

Label noise \emph{amplifies} the peak (used in the headline experiments at 15\% noise) but does not cause it. Paper Figs.\ 4a and 7 show the peak (or a plateau that develops into a peak with added noise) even at 0\% label noise. The deeper cause is model mis-specification: in any setting where the model class cannot perfectly express the conditional $\D(\bm{y}\mid\bm{x})$, fitting the train set forces the model to memorise systematic discrepancy.

\subsection{``Parameters $=$ capacity.'' \textbf{No.}}

EMC depends on the whole training procedure $\T$. Two ResNet18s with identical parameter count can have very different EMCs depending on how long they are trained, what regularisation is used, and what augmentation is applied. The paper's Fig.\ 5 (data augmentation shifts the peak) is the cleanest single demonstration.

\subsection{``EMC is a finished theory.'' \textbf{No.}}

$\varepsilon$ is heuristic ($0.1$ in the paper). ``Sufficiently smaller / larger'' is informal. The width of the critical interval is empirical, not closed-form. There is no proven non-linear theorem of EMC for over-parameterised neural networks. The hypothesis is empirically validated across many settings, but a deductive theory matching the deep-net regime is open.

% =====================================================================
\section{Q\&A reasoning templates}
% =====================================================================

For each likely question, the template is: (a) restate the question precisely, (b) work from EMC/Hyp.~1 directly rather than memorised facts, (c) acknowledge uncertainty cleanly.

\subsection{``Doesn't this just say `use early stopping'?''}

\textit{Restate.} Are you arguing that the only practical takeaway is to early-stop?

\textit{Reason.} Early stopping helps in the narrow band where the model would otherwise pass through the epoch-wise peak (Chain 2). For models much smaller than the threshold, training to completion is the right move. For models much larger than the threshold, training to completion is \emph{better} than early-stopping; that is the second descent. So the actual takeaway is regime-dependent: early stopping is one of several escape routes (Chains 4, 5).

\subsection{``How do you compute EMC in practice?''}

\textit{Restate.} How would I actually measure EMC for a given $\T$ on a given $\D$?

\textit{Reason.} You sweep $n$ for a fixed $\T$ on samples drawn from $\D$, train to completion, measure mean train error, and find the largest $n$ for which mean train error $\le\varepsilon$. This is what the paper does implicitly when they sweep widths and epochs and look at where train error first hits zero. The value of EMC depends on $\varepsilon$, which is heuristic; the paper uses $0.1$. \kw{There is no closed form.}

\subsection{``Why can the peak appear without label noise?''}

\textit{Restate.} If label noise amplifies the peak, why doesn't no-noise give no-peak?

\textit{Reason.} Chain 7. Model mis-specification is enough on its own. Real-world distributions $\D(\bm{y}\mid\bm{x})$ are rarely captured exactly by the model class $\F$, so fitting train labels exactly forces the model to memorise systematic discrepancy, which is the mechanism that label noise dramatises. Paper Figs.\ 4a (CIFAR-100, 0\% noise, peak present) and 7 (5-layer CNN, 0\% noise, peak present) confirm.

\subsection{``Is this just an artifact of CIFAR-10 / ResNet18?''}

\textit{Restate.} Could the result be specific to your headline experimental setup?

\textit{Reason.} No, the paper sweeps the cross-product:
\begin{itemize}[leftmargin=1.2em]
\item \textbf{Datasets:} CIFAR-10, CIFAR-100, IWSLT'14 De--En, WMT'14 En--Fr.
\item \textbf{Architectures:} ResNet18 (varying width), 5-layer CNN, 6-layer Transformer (varying $d_{\text{model}}$).
\item \textbf{Optimisers:} Adam, SGD with $\sqrt{T}$-schedule.
\item \textbf{Regularisation:} with and without data augmentation, with and without label smoothing, varied label noise levels.
\end{itemize}
The peak appears in essentially every cell of this matrix.

\subsection{``Why isn't this a proof?''}

\textit{Restate.} You called it a hypothesis. What is missing for it to be a theorem?

\textit{Reason.} Three missing pieces.
\begin{itemize}[leftmargin=1.2em]
\item A principled choice of $\varepsilon$.
\item A formal characterisation of ``sufficiently smaller / larger'', i.e.\ a function $\delta(\D,\T)$ that bounds the critical interval.
\item A non-linear analogue of Mei \& Montanari's closed-form for deep networks.
\end{itemize}
None of these are known. The paper validates the hypothesis empirically across an enormous experimental matrix; this is a strong scientific claim but not a deductive one.

% =====================================================================
\section{Candidate poster additions after group review}
% =====================================================================

Do not add all of these automatically. The current poster is already dense. These are the highest-value additions if the group decides that the poster needs more explicit critical thinking or implications.

\begin{enumerate}[leftmargin=1.4em]
\item \textbf{More data has two effects.} Add one sentence to the practical implications block: ``More data gives information, but it also changes where the procedure sits relative to the interpolation threshold.'' This is the strongest candidate because it explains sample-wise double descent without sounding like ``data is bad.''
\item \textbf{Noise amplifies but does not fully explain the peak.} Add one sentence to the limitations or critique block: ``Label noise raises the peak, but misspecification and optimisation also matter.'' This protects us from overclaiming.
\item \textbf{Definition / identity / hypothesis labels.} Add tiny labels near the mathematical block: risk decomposition = identity; EMC = definition; generalised double descent = hypothesis. This would show mathematical discipline with very little space cost.
\item \textbf{Practical risk-zone message.} If one implication sentence is needed: ``The dangerous region is not simply `large model'; it is $\mathrm{EMC}\approx n$, where validation, augmentation, regularisation, ensembling, or scaling away from the threshold matter most.''
\end{enumerate}

% =====================================================================
\section{Reading order for group rehearsal}
% =====================================================================

For each group member, in this order:
\begin{enumerate}[leftmargin=1.4em]
\item Read \S1 (equation index) cover-to-cover. Re-derive each equation on a whiteboard.
\item Read \S2 (reference index). Be able to say one sentence about each cited paper.
\item Pick three of the eight chains in \S3 and walk through them aloud without notes.
\item Read \S4 (anti-claims) and \S5 (Q\&A templates) before each rehearsal.
\end{enumerate}

\noindent\textbf{Hard constraint.} If pressed in Q\&A and unsure, say so. The rubric explicitly rewards honest uncertainty (point 4b: ``directly owning up not knowing the answer''). Never invent a theorem, paper, or experimental number on the fly.

\end{document}
